\chapter{Abertura Multifaixa e Polarimetria}

A abertura é deliberadamente híbrida, e não uniforme em todo o espectro. Uma Yagi VHF externa de polarização única e uma Yagi UHF externa de polarização única compartilham um boom alinhado à normal externa da face: a antena VHF usa elementos maiores e a antena UHF usa elementos menores. Suas direções transversais são mutuamente ortogonais a 45° e 135° na base tangente local. Essas orientações mecânicas fornecem diversidade entre faixas diferentes, não duas componentes de polarização de um mesmo sinal. Antenas ou tiles menores de L/S/C, X/Ku e K/Ka ocupam as camadas da face e seus suportes locais. HF utiliza loops, dipolos cruzados e diversidade de modos da plataforma. Cada família de antenas possui filtragem, conversão e amostragem próprias, enquanto a face preserva uma interface comum de sincronismo, calibração e fusão.

O primeiro demonstrador não declara cobertura espectral contínua. Os intervalos de 323--470 MHz e 860--960 MHz permanecem deliberadamente sem atribuição. Uma abertura ou array aeronáutico dedicado de 960--1215 MHz fornece dois caminhos simultâneos e independentes: preselector UAT em 978 MHz, limitador/LNA, ADC coerente e canal FPGA; e preselector 1090ES em 1090 MHz, limitador/LNA, ADC coerente e canal FPGA. Nenhum desses caminhos é atribuído à Yagi UHF de 470--860 MHz, e a abertura aeronáutica não implica monitoramento contínuo em toda a sua faixa de sintonia.

A síntese polarimétrica somente é permitida em um módulo da mesma faixa que forneça duas portas ortogonais, simultâneas e coerentes na mesma frequência (\texttt{POL-001}). Em uma base local $x$--$y$:
\[
E_{\mathrm{RHCP}}=\frac{E_x-jE_y}{\sqrt{2}},\qquad E_{\mathrm{LHCP}}=\frac{E_x+jE_y}{\sqrt{2}}.
\]
As equações não combinam as Yagis VHF e UHF. O primeiro demonstrador VHF/UHF trata essas antenas como canais independentes de polarização única (\texttt{POL-002}); a polarimetria completa exige um módulo posterior dual-porta ou cruzado na mesma faixa. Em uma superfície curva, bases locais válidas devem ser rotacionadas para um sistema comum antes da fusão. A calibração polarimétrica depende de frequência, ângulo, temperatura, equilíbrio de amplitude, fase e acoplamento mútuo.

Um mastro futuro separado de quatro canais, distinto do par VHF/UHF combinado ilustrado, implementa VHF-X/VHF-Y a 0°/90° e UHF-A/UHF-B a 45°/135°, com quatro booms e cadeias RF eletricamente independentes (\texttt{POL-004}--\texttt{POL-006}). Para cada faixa $b$, a calibração segue:
\[
\widehat{\mathbf{E}}_b(f,\theta,\phi,T)=
\mathbf{C}_b^{-1}(f,\theta,\phi,T)
\left[\mathbf{v}_b-\widehat{\mathbf{n}}_b\right],
\]
onde $\mathbf{C}_b$ é a matriz complexa de Jones $2\times2$ medida, contendo diagramas embarcados, desequilíbrio de ganho/fase, atrasos, acoplamento mútuo, mastro e plataforma. Uma pseudoinversa regularizada substitui a inversa direta quando o condicionamento for ruim. A estimativa UHF é então rotacionada por $\psi=45^\circ$ para a base global comum antes da formação dos produtos de polarização.

\section{Solução de engenharia}
O conjunto de Yagis de faixas cruzadas é uma solução proposta de integração mecânica. A Yagi maior cobre VHF e a menor cobre UHF; cada uma permanece casada, digitalizada e calibrada de forma independente. Suas saídas podem ser comparadas ou fundidas como observações de faixas diferentes, mas nenhuma operação de fase relativa produz RHCP, LHCP, Jones ou Stokes a partir desse par. Não se reivindicam novidade ou desempenho antes da busca de anterioridade e da medição.

Para ampliar a cobertura de descoberta em VHF, uma LPDA de 11 elementos para 72--320 MHz, com elemento máximo de 2 m, fica registrada como candidata de simulação (\texttt{ANT-001}--\texttt{ANT-005}). Ela ainda não substitui a Yagi VHF ilustrada: o boom de aproximadamente 2 m, a folga nas bordas, a carga de vento, a transição de alimentação, o ganho e a interação com a antena UHF exigem validação.

\begin{figure}[H]
\centering
\includegraphics[width=0.95\textwidth]{pt-BR/fig04_particionamento_espectro.png}
\caption{Particionamento espectral em escala horizontal logarítmica. Cada barra colorida representa uma família de faixa proposta; os intervalos em branco entre os caminhos VHF, UHF e aeronáutico dedicado são lacunas deliberadas do primeiro demonstrador.}
\label{fig:fig04_particionamento_espectro}
\end{figure}

\begin{figure}[H]
\centering
\includegraphics[width=0.95\textwidth]{pt-BR/fig05_polarimetria.png}
\caption{Fluxo superior: canais VHF e UHF digitalizados independentemente terminam em um resultado polarimétrico entre faixas invalidado pela cruz vermelha. Fluxo inferior: duas portas ortogonais coerentes na mesma faixa alimentam produtos calibrados de Jones, Stokes ou polarização circular.}
\label{fig:fig05_polarimetria}
\end{figure}
